\section{Results}
\subsection{Validation Results to Date}
\label{sec:validation_results}

This subsection reports evidence against the acceptance criteria of
Table~\ref{tab:requirements}, in the order the stages are entered. Where no
result exists yet, the criterion is stated so the gap is explicit rather than
silent. Appendix~\ref{app:verification} carries the per-requirement
verification status that the results below close out.

\subsubsection{1-DoF Control Validation}
\label{sec:test_1dof}

The single-axis rig is the protected experimental item of
Section~\ref{sec:schedule}, and its balance result is the evidence carried into
this document. Gate G3 requires the rig to balance for at least 60~s unassisted.

Beyond the duration criterion, three quantities are recorded because they
determine whether the result transfers to the integrated cube rather than merely
holding on the rig. Recovery from a defined impulse establishes that the balance
is a stabilised equilibrium and not a fortunate initial condition. Steady-state
wheel-speed drift measures momentum accumulation: a controller that balances
while monotonically spinning up is one that will saturate, and the drift rate
predicts when. Control effort is compared against the continuous and peak current
limits of the driver given in Appendix~\ref{app:bom}, since a balance sustained
only by exceeding the continuous rating is not a balance that survives to the
demonstration.

\paragraph{Balancing performance examples.} Table~\ref{tab:balance_examples}
reports representative runs against Gate~G3, each initialised from rest at the
upright equilibrium with no applied disturbance.

\begin{table}[htbp]
\centering
\caption[1-DoF unassisted balance examples]{Representative unassisted balance
runs on the single-axis rig, against the Gate~G3 threshold of
\qty{60}{\second}.}
\label{tab:balance_examples}
\begin{tabular}{lccc}
\toprule
Run & Balance duration & Steady-state wheel-speed drift & Peak driver current \\
\midrule
1 & $> \qty{60}{\second}$ & -- & \qty{0.36}{\ampere} \\
2 & $> \qty{60}{\second}$ & -- & \qty{0.38}{\ampere} \\
3 & $> \qty{60}{\second}$ & -- & \qty{0.40}{\ampere} \\
\bottomrule
\end{tabular}
\begin{minipage}{\linewidth}
\vspace{4pt}
\footnotesize
Note: No constant drift was appreciated during static zero-position balancing. When the center of mass was uncalibrated, small oscillations were observed instead of steady drift due to torque saturation.
\end{minipage}
\end{table}

\paragraph{Recovery-angle disturbance tests.} Beyond the zero-disturbance case,
the rig was displaced by hand to a defined initial angle, measured from the
support, and released to test whether the controller recovers to upright
rather than merely holding a small perturbation. Table~\ref{tab:recovery_angle}
reports the three angles used to characterise this margin.

\begin{table}[htbp]
\centering
\caption[Recovery from a defined initial angle]{Recovery outcome versus the
initial displacement angle, measured from the support.}
\label{tab:recovery_angle}
\begin{tabular}{lccc}
\toprule
Initial angle & Recovered? & Recovery time & Peak driver current \\
\midrule
\qty{7}{\degree}  & Yes & \qty{1.2}{\second} & \qty{0.86}{\ampere} \\
\qty{8}{\degree}  & Yes & \qty{1.5}{\second} & \qty{0.98}{\ampere} \\
\qty{12}{\degree} & Yes & \qty{2.1}{\second} & \qty{1.34}{\ampere} \\
\bottomrule
\end{tabular}
\end{table}

The largest initial angle tested to date, again measured from the support, is
\qty{24}{\degree}; recovery from this angle was also satisfactory, so the
demonstrated disturbance-rejection envelope currently extends at least this
far beyond the three angles of Table~\ref{tab:recovery_angle}. This angle represents
the physical hard stop of the experimental rig support as well as the maximum tilt
evaluated in simulation, so requirement PER-03 (recovery from
\qty{12}{\degree} or more) is satisfied with margin, and the demonstrated
envelope is bounded by the rig rather than by the controller.

\subsubsection{3-DoF Control Validation}
\label{sec:test_3dof}

For the 3D model, control-loop validation to date has achieved edge balance on
the final prototype using the 1-DoF control law; corner balance has not yet
been demonstrated.
